\section{チェッカーボードモードの代数的証明}
\label{app:checkerboard_mode}
% ═══════════════════════════════════════════════════════════════════════════════

格子間隔 $h$ の一様格子上でチェッカーボードモード
\begin{equation}
  p_j = P_0(-1)^j
\end{equation}
を考える．中心差分による圧力勾配は：
\[
  \left.\frac{dp}{dx}\right|_j \approx \frac{p_{j+1} - p_{j-1}}{2h}
  = \frac{P_0(-1)^{j+1} - P_0(-1)^{j-1}}{2h}
  = \frac{-P_0(-1)^j - (-P_0(-1)^j)}{2h} = 0
\]
すなわち，$p_j = P_0(-1)^j$ は（定数圧力を除いて）あらゆる物理的な圧力から区別できない非物理的なモードであり，
中心差分では\textbf{零勾配を与える（ゴーストモード）}となる．

このモードが圧力場に混入すると，速度場を修正しないまま圧力のみ発散する「チェッカーボード不安定」を引き起こす
（対策は §\ref{sec:checkerboard_solution} および §\ref{sec:balanced_force} で詳述している）．

% ═══════════════════════════════════════════════════════════════════════════════
